\section{Trabalhos Correlatos}

A identificação de patologias cardíacas envolve múltiplas etapas, e um histórico cardíaco detalhado é essencial nesse processo \cite{malik2023cardiac}. Diversos estimadores de risco foram propostos para identificar níveis de risco de doenças cardíacas, como o Framingham Risk Score \cite{hemann2007framingham} e o QRISK \cite{hippisley2007derivation}, no entanto, esses estimadores apresentam limitações ao não considerarem as variabilidades individuais dos pacientes \cite{talha_major_2024}.

A inteligência artificial tem sido amplamente utilizada para a detecção de patologias cardíacas, como em \cite{subramani_cardiovascular_2023}, onde os autores avaliaram diferentes modelos individualmente e também em conjunto, encontrando que a fusão de modelos resulta em uma performance superior em relação aos modelos individuais. Por outro lado, \cite{dahia_implementing_2023} avaliou diversos modelos de aprendizado de máquina e conseguiu atingir uma medida AUC de 87,3\% com os modelos Adaptive Boosting e Random Forest. Apesar dos resultados promissores na detecção de doenças cardíacas, algumas limitações ainda persistem.

\cite{chakraborty2017interpretability} levanta a questão sobre a interpretabilidade dos modelos de aprendizado de máquina, onde a complexidade desses modelos dificultam a interpretação dos resultados. A interpretabilidade é um fator crucial em aplicações médicas, pois os médicos precisam entender os motivos que levaram um modelo a tomar uma determinada decisão \cite{stiglic2020interpretability}. Além disso, há limitações relacionadas à disponibilidade de dados para o treinamento dos modelos \cite{fardouly2022potential}. Para superar essas limitações, diferentes abordagens têm sido propostas.

Um modelo de aprendizado por agrupamento utilizando o classificador de suporte de vetores quântico foi proposto por \cite{abdulsalam2023explainable}. O modelo, além de apresentar resultados interpretáveis, conseguiu atingir uma acurácia de 90,16\% na detecção de doenças cardíacas. Outra abordagem foi proposta por \cite{linardos2022federated}, onde os autores apresentaram um estudo de aprendizado federado com imagens de ressonância magnética cardíaca, no qual o modelo foi treinado com dados provenientes de quatro centros diferentes.